\documentclass[11pt,letterpaper]{article}
\usepackage[margin=0.8in]{geometry}
\usepackage[T1]{fontenc}
\usepackage{lmodern,amsmath,amssymb,booktabs,tabularx,url}
\usepackage[hidelinks]{hyperref}
\setlength{\parindent}{0pt}\setlength{\parskip}{5pt}
\setlength{\emergencystretch}{2em}
\providecommand{\reportbuilddatetime}{unknown}
% BEGIN BUILD METADATA
% Generated by build_report.py.
\begin{filecontents*}[overwrite]{methods_build_info.tex}
\renewcommand{\reportbuilddatetime}{2026-10-07 13:06:51 EDT}
\end{filecontents*}
% END BUILD METADATA
\InputIfFileExists{methods_build_info.tex}{}{}
\title{Sparse metric MDS in grip\\\large Euclidean landmarks, regional interactions and graph geodesics}
\author{\small\parbox{0.95\textwidth}{\centering Source: \path{~/current_projects/grip/vignettes/articles/sparse-mds-methods/methods.tex}}}
\date{Build \reportbuilddatetime}
\begin{document}\maketitle
\section{Purpose and status}
This note specifies three sparse metric multidimensional scaling (MDS) methods
and their implementation in \texttt{grip}. Two methods approximate Euclidean
pairwise distances. The third approximates shortest-path distances on a connected
neighbor graph. A change of target geometry is not an improvement in approximation
of the same target. The numerical backend is stochastic gradient descent (SGD).
No full distance matrix is constructed by the sparse preparation or fitting path.
Small dense reference calculations are permitted only in explicit validation.

\textbf{Status: implemented and tested on the synthetic examples below.}
The new observation preparation computes pivots, local interactions and regional
weights internally. It supports Euclidean and graph-geodesic targets, with
streaming component-MST repair in \texttt{dgraphs}. Uniform landmarks remain an
explicit baseline. The validation compares 288 finite-budget fits and does not
claim convergence or optimal tuning for the combined V4 data.

Let $x_i\in\mathbb R^p$, $i=1,\ldots,n$, be finite distinct input profiles, and
$y_i\in\mathbb R^k$ their embedding. Exact duplicates must be consolidated before
inverse-squared weighting. The input units define Euclidean distance
$d^E_{ij}=\|x_i-x_j\|_2$. A connected graph with positive Euclidean edge lengths
defines $d^G_{ij}$ as the minimum sum of lengths over paths from $i$ to $j$.
All computations within a variant use its chosen target metric consistently.

\begin{table}[ht]\centering\small
\caption{The three variants. Regional multipliers approximate omitted interactions;
they are distinct from inverse-squared target-distance weights.}
\begin{tabularx}{\textwidth}{lXXX}\toprule
Variant & Pivot selection & Local pairs & Target metric\\\midrule
Uniform landmarks & Uniform without replacement & None & Euclidean\\
Euclidean sparse stress & Distance-proportional; optional farthest-first & Symmetric neighbor graph & Euclidean\\
Geodesic sparse stress & Distance-proportional & Repaired neighbor graph & Graph shortest paths\\\bottomrule
\end{tabularx}\end{table}

The sparse-stress reference is Zheng et al., Section IV-C1 and Algorithm 2
\cite{zheng2018sparse}, adapting Ortmann et al. The point-cloud graph construction,
PCA threshold and example sizes here are implementation choices, not prescriptions
from that paper. The legacy \texttt{grip} learning-rate schedule is distinguished
from the published weight-dependent schedule below.

\clearpage\section{Local graph and its search coordinates}
Center the input to obtain $Z$. Ordinary PCA uses its right singular vectors
$V$ and variances $\lambda_1\geq\cdots\geq\lambda_p\geq0$. For retained-variance
target $t=0.90$, configurable, set
\[
r=\min\{a:\sum_{j=1}^a\lambda_j/\sum_j\lambda_j\geq t\},\qquad U=ZV_r.
\]
Search on rows of $U$, without whitening or column standardization. A null target
requests direct input-space search. Let $N_q(i)$ be the $q$ nearest other rows;
retain $\{i,j\}$ if $j\in N_q(i)$ or $i\in N_q(j)$. Use $q=10$ by default,
with 5 and 20 for sensitivity. Search and target geometry are separate: every
retained edge is measured using $d^E_{ij}$ from $X$, never the truncated scores.
The bundled ANN library currently runs at zero approximation tolerance, hence
exact neighbor search in the chosen search coordinates. PCA itself can change
neighbors. Zero projected distances are valid search ties, not new duplicates.

Projection preserves the fraction $\sum_{i<j}\|u_i-u_j\|^2/
\sum_{i<j}\|x_i-x_j\|^2=\sum_{j\leq r}\lambda_j/\sum_j\lambda_j$ of aggregate
squared distances. It does not bound each pair's distortion or ensure neighbor
preservation. Record the achieved variance, dimension, center, loadings and
ordered identifiers. Reuse these coordinates for comparisons across $q$.

\subsection{Connected graph without a dense distance matrix}
Only the geodesic variant requires connectivity repair. If the neighbor graph
has $c$ components, define the length between two components as the shortest
original-space Euclidean edge joining their vertices. Add a minimum spanning
tree (MST) of this implicit component graph, giving exactly $c-1$ bridges.

Use Prim's algorithm, beginning with the component containing vertex one.
Maintain, for each unvisited component, its best bridge from the visited set.
When a component enters the tree, evaluate distances from its vertices to
vertices of unvisited components and update those bridges. Select the next
bridge by length, then lexicographic endpoint order to resolve ties. Distances
are consumed immediately; no $n\times n$ or $c\times c$ table is retained.

This is an exact component MST with $O(n+c)$ auxiliary storage. The worst-case
distance work remains $O(n^2p)$: avoiding a matrix is not a promise of subquadratic
work. Record the number of evaluated cross-component pairs and added bridges.
An ANN-candidate repair may fall back to this streaming algorithm, never to a
dense distance allocation. Sparse fitting then uses all original and added graph
edges as local interactions. Bridges impose relationships between previously
separate groups; connectivity does not establish a biological manifold.

\clearpage\section{Pivots, distance rows and regional weights}
Use $h$ pivots, with $1\leq h<n$ for observation preparations.
Requested counts are capped at $n-1$; the all-pivot limit is tested separately. The uniform baseline samples indices without
replacement and retains every unordered pair touching a pivot once. Its endpoint
multipliers are one. It has no regional weights but does have inverse-squared
weights: its stress is $\sum_{\mathcal F}(\|y_i-y_j\|-d_{ij})^2/d_{ij}^2$,
where $\mathcal F$ is the retained pair set.

For regional sparse stress, choose the first pivot uniformly. Given selected
pivots $P$, maintain $r_i=\min_{v\in P}d_{iv}$. Select the next pivot with
probability $r_i/\sum_j r_j$, using distance rather than squared distance.
The Euclidean alternative selects an index maximizing $r_i$, resolving ties by
index. In either case compute only one pivot-to-all distance row per selected
pivot. Euclidean rows use direct coordinate differences; graph rows use Dijkstra's
algorithm. Save the rows for reuse, region construction and audit.

The region $R_v$ of pivot $v$ contains profiles closest to it, with ties assigned
to the earliest selected pivot. For a nonlocal pivot--profile pair, define
\[
s_{iv}=\#\{j\in R_v:d_{vj}\leq d_{vi}/2\},\qquad w'_{iv}=s_{iv}/d_{iv}^2.
\]
This weight controls movement of endpoint $i$. The reverse weight is zero unless
$i$ is also a pivot, in which case compute its own regional count. Sort each
region's saved distances to its owner and count values below the inclusive
half-distance threshold by binary search. All profiles count once; input-frequency
multiplicities do not enter this construction.

For each local pair, override both endpoint multipliers with one. Do not add
local and regional contributions for overlapping pairs. Enumerate each unordered
pair once and omit self-pairs. The number of constraints is at most
$hn-h(h+1)/2+m$, where $m$ is the local edge count. Euclidean distance preparation
requires $O(hnp)$ arithmetic and $O(hn+m)$ storage. Graph preparation requires
$h$ single-source shortest-path searches, not all-pairs paths.

If $h=n$, the regional counts equal one and the construction reduces to full
inverse-squared stress, provided graph edge lengths are themselves shortest-path
distances. Euclidean edge lengths satisfy that condition by the triangle
inequality. Arbitrary weighted input graphs can contain longer redundant edges;
the legacy kernel preserves such edge-length constraints, which must not be
mistaken for full shortest-path stress.

\clearpage\section{SGD updates and reported scores}
For a pair $i<j$, let $a_{ij},b_{ij}$ be its two endpoint multipliers and let
$w_i=a_{ij}/d_{ij}^2$, $w_j=b_{ij}/d_{ij}^2$. Write $\ell=\|y_i-y_j\|$.
For positive $\ell$, the half-distance correction and update are
\[
v=\frac{\ell-d_{ij}}{2\ell}(y_i-y_j),\quad
 y_i\leftarrow y_i-\min(\eta w_i,1)v,\quad
 y_j\leftarrow y_j+\min(\eta w_j,1)v.
\]
A zero-weight endpoint remains fixed in that update. Use the existing seeded
collision handling when $\ell=0$. Each epoch shuffles and visits the retained
unordered pairs once. The published fixed-budget exponential schedule uses
$\eta_{\max}=1/w_{\min}$ and $\eta_{\min}=\varepsilon/w_{\max}$, where extrema
include only positive endpoint weights and $\varepsilon=0.1$. Interpolate
linearly between their logarithms over the epoch indices, including endpoints;
a one-epoch run uses $\eta_{\max}$. Normalize targets by their root-mean-square
value before fitting, and use weights in those same normalized units.

Expose this as the named published schedule, retaining the existing hybrid
schedule for backward compatibility. Record the effective rates, seeds, output
dimension, start count, epoch budget, pair count and achieved scores. Return the
terminal iterate; completion of the epoch budget is not convergence.

The scalar monitoring proxy is
\[
S(Y)=\sum_{\{i,j\}\in\mathcal F}\frac{a_{ij}+b_{ij}}{2}
 \left(\frac{\|y_i-y_j\|-d_{ij}}{d_{ij}}\right)^2.
\]
With asymmetric multipliers, its gradient is not the endpoint update above.
Report it explicitly as a proxy. Independently evaluate layouts using
\[
R(Y;\mathcal A)=100\sqrt{\frac1{|\mathcal A|}\sum_{\{i,j\}\in\mathcal A}
 (\|y_i-y_j\|/d_{ij}-1)^2}.
\]
Here $\mathcal A$ is a declared common evaluation set, $R$ is RMS relative error
in percent, and smaller is better. Evaluate against each method's own target
metric. Dense reference results are achieved finite-budget solutions, not proven
optima. Also report unscaled PCA on the same metric as a descriptive baseline.
For matched numerical comparisons, small dense reference objectives are
represented by all unordered pair constraints through the same SGD interface.
This keeps terminal-iterate selection and scaling conventions identical across
methods. No fitted rescaling or best-start selection is applied in validation.

\clearpage\section{Public interface and reproducibility}
The entry point is \texttt{prepare.sparse.mds(X, method, ...)}.
It computes pivots and weights internally and returns reusable sparse preparation
for \texttt{metric.mds(prepared=..., approximation="sparse", backend="sgd")}.
Methods are \texttt{uniform}, \texttt{euclidean} (default) and \texttt{geodesic}.
Controls include pivot count, selection rule, neighborhood size, PCA retained
variance, preparation seed and workspace allowance. Fit seeds and budgets belong
to \texttt{metric.mds()}, so multiple fits can share one preparation.
\begin{verbatim}
prep <- grip::prepare.sparse.mds(X, method="euclidean",
  n.pivots=200, q=10, neighbor.variance.target=0.90, seed=1)
fit <- grip::metric.mds(prepared=prep, approximation="sparse",
  backend="sgd", dim=3, max.iter=100, seed=1,
  sgd.control=list(scheduler="zheng", schedule.epsilon=0.1))
\end{verbatim}
These illustrative counts are configurable, not a V4 calibration result.
Use \texttt{method="uniform"} for the landmark baseline,
\texttt{pivot.selection="farthest"} for Euclidean farthest-first selection,
and \texttt{method="geodesic"} for connected-graph targets.

The preparation records input identity, algorithm choices, PCA information,
graph edges and bridge diagnostics, ordered pivots, owners, distance rows,
constraint targets and endpoint multipliers. Invalid or duplicate inputs fail
explicitly. Allocation estimates are checked before large arrays; they are
not a process-wide memory quota. Code and compiled-library fingerprints belong
in run manifests, and cache reuse requires matching data and settings.

\section{Validation design and examples}
Quadratic-form examples have latent dimension $d=2$ or $3$, coordinates
$u\in[-1,1]^d$ and observations
\[
x(u)=(u,\;0.8u^{\mathsf T}Au),\qquad
 A=I_d\ \text{(index 0)},\quad A=\operatorname{diag}(-1,1,\ldots,1)
 \ \text{(index 1)}.
\]
Index means the number of negative eigenvalues of the quadratic form. Thus
these are 2D surfaces in three coordinates and 3D hypersurfaces in four
coordinates, not merely 2D and 3D display layouts. Use 120 seeded uniform latent
samples per case, both 2D and 3D output, three seeds, 100 epochs, and 16 pivots.
Compare uniform, randomized Euclidean, farthest-first Euclidean, geodesic and
dense references for each target geometry. Vary $q=5,10,20$ for the two regional
families, keeping data, initial coordinates and seeds matched. Use a separated
cluster fixture to force connectivity repair, and a flat fixture for exact
Euclidean reconstruction checks.

Independent checks compare pivot rows, ownership, half-radius counts, local
pair override, all-pivot reduction, RNG preservation, repeatability, unit
invariance and allocation refusal. Compare streamed repair against a small dense
component-MST reference, including ties and a forced ANN fallback. Confirm that
production preparation has no full distance matrix. Small validation matrices
are deliberately constructed outside the sparse path and clearly identified.

% BEGIN VALIDATION RESULTS

\clearpage\section{Measured results}
The final experiment completed all 288 attempted fits: zero failed, zero skipped,
and zero excluded. Each method--case--output combination has three preparation
and fitting seeds, with one start per fit. Each layout was evaluated on all
7,140 unordered pairs in its own target metric. These small dense evaluations
are validation calculations, outside the production sparse path.

\subsection{Euclidean approximation}
Table~\ref{tab:euclidean} shows the median RMS relative distance error, $R$, in
percent. The randomized regional method at $q=10$ improved on uniform landmarks
in 23 of 24 paired comparisons (one worse, no ties); farthest-first improved in
24 of 24. The much larger errors of the uniform baseline on the 3D forms embedded
in two dimensions illustrate a limitation of using landmark interactions alone.
No method wins universally from these results: the experiment uses four fixed
point clouds and only three seeds, without tuning pivot counts or epoch budgets.

\begin{table}[ht]\centering\small
\caption{Euclidean targets: median $R$ (percent; lower is better), over three
seeds. $d$ is latent dimension, $i$ is quadratic-form index, and $k$ is output
dimension. PCA is the deterministic unscaled projection. Dense is the all-pair
SGD reference; Uniform uses 16 random landmarks; Random and Farthest use 16
regional pivots and $q=10$. Values below 0.001\% are shown as such. All fits
are included; there are no missing values.}\label{tab:euclidean}
\begin{tabular}{rrr rrrrr}\toprule
$d$ & $i$ & $k$ & PCA & Dense & Uniform & Random & Farthest\\\midrule
2 & 0 & 2 & 14.15 & 8.04 & 9.26 & 8.21 & 8.22 \\
2 & 0 & 3 & $<0.001$ & $<0.001$ & 0.66 & $<0.001$ & $<0.001$ \\
2 & 1 & 2 & 15.25 & 7.32 & 8.38 & 7.49 & 7.43 \\
2 & 1 & 3 & $<0.001$ & $<0.001$ & 2.84 & $<0.001$ & 0.001 \\
3 & 0 & 2 & 32.20 & 24.88 & 37.62 & 26.85 & 26.05 \\
3 & 0 & 3 & 12.22 & 8.29 & 10.99 & 8.72 & 8.82 \\
3 & 1 & 2 & 32.65 & 24.07 & 55.60 & 26.76 & 25.68 \\
3 & 1 & 3 & 12.06 & 6.99 & 9.10 & 7.99 & 7.29 \\
\bottomrule\end{tabular}\end{table}

Against the matched dense SGD reference, the randomized regional method's median
paired excess error was 0.30 percentage points, ranging from approximately zero
to 2.97. Farthest-first had median excess 0.24 points and maximum 4.14. These are
pooled descriptions of 24 paired runs; Table~\ref{tab:euclidean} and the linked
figures retain the individual cases and output dimensions. No sparse run beat
its dense reference on this evaluation. The reference is a finite-budget fit,
not a demonstrated optimum or a lower bound.

The index-0 and index-1 2D surfaces are represented in three ambient coordinates;
three output dimensions therefore admit exact Euclidean reconstruction. Dense
SGD and PCA reached numerical precision there. Regional fits approached that
limit but were not identically exact. In contrast, the 3D forms occupy four
ambient coordinates, so their three-coordinate Euclidean embeddings remain
approximate. These observations distinguish latent dimension from the number
of Euclidean coordinates required to reproduce all chord distances.

\clearpage\subsection{Graph geometry and neighborhood sensitivity}
Graph-geodesic stress has a different target. In Table~\ref{tab:geodesic}, both
columns use shortest paths on the same $q=10$ graph. The regional method's median
paired excess error over its dense reference was 0.64 percentage points, with
range 0.21--4.01. Its nonzero errors for 2D surfaces embedded in three dimensions
are compatible with the altered target geometry; they are not evidence that
the Euclidean implementation lost an exact reconstruction.

\begin{table}[ht]\centering
\caption{Graph-geodesic targets at $q=10$: median $R$ (percent) over three seeds.
$d$, $i$ and $k$ have the meanings in Table~\ref{tab:euclidean}. Dense fits all
7,140 pairs; Regional uses 16 internally selected pivots and local edges.
Each row uses the same graph target in both columns.}\label{tab:geodesic}
\begin{tabular}{rrr rr}\toprule
$d$ & $i$ & $k$ & Dense & Regional\\\midrule
2 & 0 & 2 & 7.73 & 7.99 \\
2 & 0 & 3 & 5.96 & 6.24 \\
2 & 1 & 2 & 6.76 & 7.28 \\
2 & 1 & 3 & 5.00 & 5.54 \\
3 & 0 & 2 & 24.57 & 26.13 \\
3 & 0 & 3 & 9.57 & 10.24 \\
3 & 1 & 2 & 23.39 & 24.96 \\
3 & 1 & 3 & 8.91 & 9.70 \\
\bottomrule\end{tabular}\end{table}

Changing $q$ from 5 to 10 to 20 gave pooled median paired excess errors of
0.33, 0.30 and 0.13 percentage points for Euclidean targets, and 1.05, 0.64 and
0.30 for graph targets. The graph reference changes with $q$; those latter
numbers compare approximation error within each graph, not the biological or
geometric merits of different graphs. Improvement is not asserted for every seed.
The \href{../../../output/sparse-mds-methods/neighborhood-sensitivity.pdf}
{neighborhood-sensitivity supplement (Figures 5--8)} shows all four quadratic
forms, both output dimensions, and all three seeds at every $q$.
The \href{../../../output/sparse-mds-methods/error-comparison.pdf}
{primary comparison supplement (Figures 1--4)} shows seed-level errors at $q=10$.
Gray points are individual seeds; red markers and segments show medians and
observed ranges. These figures are descriptive. Three seeds and one sampled
cloud per form provide insufficient replication for population uncertainty
claims; ranges are not confidence intervals.

All quadratic-form graphs were connected before repair. All retained their
full three or four PCA coordinates at the 90\% threshold, with complete neighbor
agreement. Consequently these fits do not validate PCA truncation or repair
performance. Separate unit fixtures cover a 50-coordinate lifted cloud with
actual reduction, projected coincidences with positive original distances,
and three separated clusters requiring bridges. They establish implementation
properties, not the adequacy of 90\% retained variance for microbiome profiles.

\clearpage\section{Implementation verification and limits}
The complete \texttt{grip} test suite passed 6,665 checks, with no failures,
warnings or skips. The 112 targeted \texttt{dgraphs} checks also passed without failures, warnings
or skips. These regression tests cover exact
streamed component-MST cost and bridge count against a small dense reference,
tie reproducibility, fallback repair, PCA search coordinates, and original-space
edge lengths. Independent \texttt{grip} checks compare saved distance rows,
region counts and targets with direct calculations; test the all-pivot limit,
2D/3D fits, published learning rates, scaling, RNG preservation and memory refusal.

The first synthetic dispatch stopped before fitting because the new function
was not yet exported. Regenerating package documentation and exports resolved
that setup failure. A later edge-case test exposed rejection of coincident PCA
scores; the search-coordinate validator was corrected to retain those distinct
original observations. The final 288-fit run was repeated after that correction,
and its source and compiled-library fingerprints were unchanged throughout.
Earlier attempts and their logs remain available.

The final R process took 41.12 seconds of serial wall time and had maximum
resident memory 443,318,272 bytes (422.8 MiB). This includes package loading,
preparation, dense validation references, fitting, scoring and result saving;
it excludes native compilation and report rendering. It is not a sparse-only
memory benchmark. Fit-only timings and preparation timings are recorded
separately in the saved table. No large-$n$ runtime or memory-scaling claim is
made. The implicit-component repair avoids quadratic storage but can still
perform quadratic distance work.

Package documentation and API inventories were regenerated and checked.
The package-wide \texttt{make check-fast} command stopped at its repository
hygiene gate because an existing, unrelated serialized dataset contains a
personal filesystem path. That object was preserved. Thus the new tests pass,
but this work does not establish a clean package build or CRAN readiness.

\section{Reproduction and next use}
The results were generated at \textbf{2026-10-07 12:47:48 EDT}. This report was rebuilt
at \reportbuilddatetime\ from saved artifacts; rendering does not rerun fits.
The authoritative result bundle is
\href{../../../output/sparse-mds-methods/validation_2026-10-07_final/validation.rds}
{validation\_2026-10-07\_final/validation.rds}, with
\href{../../../output/sparse-mds-methods/validation_2026-10-07_final/fits.csv}{all 288 scores},
\href{../../../output/sparse-mds-methods/validation_2026-10-07_final/source_manifest.csv}
{source and library fingerprints}, saved graphs, preparations, coordinates and
session information. The adjacent \href{README.md}{reproduction record}
identifies commands and logs; \href{build_manifest.json}{the build manifest}
fingerprints the report and evidence.

The implementation now supports the requested three variants without constructing
a full distance matrix during sparse preparation or fitting. For the planned
combined V4 analysis, use Euclidean regional preparation as agreed, retain
$q=10$ with 5/20 sensitivity, and measure actual PCA neighbor retention on the
unique abundance profiles. These synthetic tests do not select an adequate
pivot count, epoch budget or PCA threshold for that dataset. No combined V4
embedding was computed in this methods-validation run.

% END VALIDATION RESULTS

\begin{thebibliography}{1}
\bibitem{zheng2018sparse} J. X. Zheng, S. Pawar and D. F. M. Goodman.
\emph{Graph Drawing by Stochastic Gradient Descent}. arXiv:1710.04626v3,
28 June 2018. \url{https://arxiv.org/abs/1710.04626v3}.
\end{thebibliography}
\end{document}
